\section{Two-Step QAOA as a Compilation Pass}
\label{sec:software}

The construction of Sec.~\ref{sec:theory} needs nothing beyond what a user already writes for a
penalty-based QUBO. We therefore present Two-Step QAOA as a compilation pass that turns a
standard penalty formulation into a constraint-preserving circuit without any penalty tuning
for the hardened constraints.

\paragraph{Input.} A cost QUBO $C(x)$ and a list of equality constraints $\sum_{i\in g}x_i=k_g$,
exactly as passed to a penalty builder such as \texttt{add\_onehot} / \texttt{add\_sum\_eq}.
If only a penalised QUBO matrix is available, the supports $g$ and weights $A_g$ can be suggested
by the detector of Sec.~\ref{sec:detect}, but the whole proposed decomposition (supports, $k_g$ and $A_g$) must be confirmed by the user (Remark in Sec.~\ref{sec:theory}).

\paragraph{Pass.} (The experiments of Sec.~\ref{sec:results} implement steps 2--4 for the groups chosen by the
formulation code; the greedy selection of step 1 is implemented in the detector but was not used in the experiments.)
\begin{enumerate}
\item \emph{Select hard groups.} Choose a pairwise-disjoint subset $\mathcal G$ of the constraints
  (we take groups greedily in order of decreasing number of qubits). Overlapping constraints that are not selected
  (e.g.\ the columns of an assignment matrix when its rows are hardened) stay in the cost as penalties.
\item \emph{Stage 1 from a table.} For each $g\in\mathcal G$ look up stage-1 angles for $(|g|,k_g)$
  in a precomputed table (Table~\ref{tab:stage1}); no per-problem optimisation is performed.
  Groups are independent, so different groups may use different depths and angles.
\item \emph{Stage 2.} Build the cost layer from $C$ plus the penalties of the non-hardened constraints,
  and the XY mixer of Eq.~\eqref{eq:xyproduct} on the hardened groups ($X$ mixer elsewhere).
  Train the stage-2 angles on the expectation post-selected on $\Pi$.
\item \emph{Sample and post-select.} Discard samples that violate a hardened constraint; return the best
  remaining sample. By Proposition~1 the returned answer always satisfies the hardened constraints.
\end{enumerate}

\paragraph{What the user no longer does.} (i) Choose or search a penalty weight for the hardened
constraints; (ii) design a Dicke-state circuit; (iii) optimise stage-1 parameters per problem.
What remains: a penalty weight for constraints that cannot be hardened (overlapping or inequality
constraints), and the usual stage-2 training.

\paragraph{The stage-1 table.} Because stage~1 stays in the symmetric subspace, the table is computed
with $(n+1)$-dimensional matrices. We optimise layer by layer (each depth is initialised from the
previous solution by interpolation and by appending a zero layer, plus random starts over the full
periods $\gamma\in[-\pi,\pi]$, $\beta\in[-\pi/2,\pi/2]$), so the tabulated $P_k$ never decreases with depth.
